%=============================================================================!
% 6.6  CONSTRAINTS AND LAGRANGE MULTIPLIERS                                   !
%=============================================================================!
\section{Constraints and Lagrange multipliers}
\label{sec:la-multipliers}

Choosing coordinates that obey a constraint makes its force disappear,
which is what we want when only the motion matters. Sometimes the force
itself is wanted, such as the tension of a pendulum's rod, and sometimes no
convenient coordinates exist, as for the bonds of a large molecule held at
fixed lengths. Then we keep the Cartesian coordinates and let the
constraint enter through a new unknown, a Lagrange multiplier. This
section introduces the tool, uses it on the pendulum, and compares the
result with Newton's laws.

\begin{toolbox}[Lagrange multipliers]
Suppose we want the largest or smallest value of $f(x, y)$ among the
points on the curve $\chi(x, y) = 0$, written with the Greek chi since $g$
is the acceleration of free fall. Move along the curve, in the direction
of a short arrow $\vb{u}$ along it. By the chain rule along a path
(\cref{sec:en-gradient}), $f$ changes at the rate $\grad f\cdot\vb{u}$,
and $\chi$ at the rate $\grad\chi\cdot\vb{u} = 0$, since $\chi$ stays zero
on the curve. At a largest or smallest value of $f$ on the curve its rate
of change along the curve is zero too, $\grad f\cdot\vb{u} = 0$. In the
plane there is only one direction at right angles to $\vb{u}$, so $\grad
f$ and $\grad\chi$, the latter not zero, point along the same line there:
\begin{equation*}
   \grad f = \lambda\,\grad\chi , \qquad \chi = 0 ,
\end{equation*}
for some number $\lambda$, called the Lagrange multiplier: three equations
for the three unknowns $x$, $y$ and $\lambda$. For example, the rectangle
of largest area $f = xy$ with a perimeter of \SI{4}{\metre}, $\chi = 2x +
2y - 4 = 0$, has $(y, x) = \lambda\,(2, 2)$, so $x = y$, and the perimeter
gives $x = y = \SI{1}{\metre}$: a square. With more variables the same
argument, applied to every arrow $\vb{u}$ along the surface $\chi = 0$,
again puts $\grad f$ along $\grad\chi$; with several constraints $\chi_c =
0$, one multiplier $\lambda_c$ each, it becomes $\grad f =
\sum_c\lambda_c\grad\chi_c$, which we take on trust.
\end{toolbox}

Here $\lambda$ is a multiplier, not the eigenvalue of \cref{sec:os-modes}.

\subsection*{The pendulum in $x$ and $y$}

Describe the bob by its Cartesian position $\vb{r} = (x, y)$, with the
pivot at the origin and $y$ up, and impose the rod as the constraint $\chi
= \frac12(x^2 + y^2 - L^2) = 0$. Its gradient is $\grad\chi = (x, y) =
\vb{r}$, along the rod and at right angles to the circle $\chi = 0$, as a
gradient is to its contour (\cref{sec:en-gradient}). The rod does no work
(\cref{sec:la-coordinates}), so its force too is at right angles to the
circle, along $\grad\chi$: it is $\lambda\,\grad\chi = \lambda\vb{r}$ for
some multiplier $\lambda$ that may change with time, and Newton's second
law reads
\begin{equation*}
   m\ddot{\vb{r}} = m\vb{g} + \lambda\vb{r} ,
\end{equation*}
with $\vb{g} = (0, -g)$ the acceleration of free fall.

Hamilton's principle gives the same equation, and shows where the
multiplier comes from. Here $\Lag = \frac12 m(\dot{x}^2 + \dot{y}^2) -
mgy$. Changing the path by $\zeta\boldsymbol{\eta}(t)$, with
$\boldsymbol{\eta} = (\eta_x, \eta_y)$ zero at both ends, the steps of
\cref{sec:la-euler}, once for $x$ and once for $y$, change the action by
$\zeta\int\vb{B}\cdot\boldsymbol{\eta}\,\dd t$ plus terms of order
$\zeta^2$, where
\begin{equation*}
   \vb{B} = \Bigl(\frac{\partial\Lag}{\partial x}
   - \frac{\dd}{\dd t}\frac{\partial\Lag}{\partial\dot{x}},\
   \frac{\partial\Lag}{\partial y}
   - \frac{\dd}{\dd t}\frac{\partial\Lag}{\partial\dot{y}}\Bigr)
   = \bigl(-m\ddot{x},\ -mg - m\ddot{y}\bigr)
   = m\vb{g} - m\ddot{\vb{r}} .
\end{equation*}
Among all paths this would give $\vb{B} = \vb{0}$, the bob falling freely.
Only paths that keep the bob on its circle are allowed, so at each instant
$\boldsymbol{\eta}$ must point along the circle, $\boldsymbol{\eta} =
\eta(t)\,\vb{e}_\theta$ with $\vb{e}_\theta$ the arrow of length one along
the circle at the bob (\cref{eq:ro-polar}); such a change alters the length only by terms in $\zeta^2$, which
do not touch the term in $\zeta$. The change of action is then
$\zeta\int(\vb{B}\cdot\vb{e}_\theta)\,\eta\,\dd t$, and the bump argument
of \cref{sec:la-euler} gives $\vb{B}\cdot\vb{e}_\theta = 0$ at every
instant:
$\vb{B}$ has no part along the circle. In the plane only the direction of
$\vb{r}$ is left, as in the toolbox, so $\vb{B} = -\lambda\vb{r}$ for some
number $\lambda$ at each instant, the minus sign only a choice of name,
and $m\vb{g} - m\ddot{\vb{r}} = -\lambda\vb{r}$ is the second law above.
The constraint leaves the principle silent about one force, along
$\grad\chi$, and that force is the rod's. Equivalently, these are the
Euler--Lagrange equations of $\Lag + \lambda\chi$ in $x$ and $y$ treated
as free, the usual way of writing a multiplier, with one $\lambda$ for
each instant.

The multiplier is fixed by the constraint itself. Differentiating
$\vb{r}\cdot\vb{r} = x^2 + y^2 = L^2$ twice with respect to time, by the
product rule on each component as in \cref{sec:mo-circle},
\begin{align*}
   2\,\vb{r}\cdot\dot{\vb{r}} &= 0 ,\\
   \dot{\vb{r}}\cdot\dot{\vb{r}} + \vb{r}\cdot\ddot{\vb{r}} &= 0
   && \text{differentiating again and halving,}\\
   \lvert\vb{v}\rvert^2 + \vb{r}\cdot\vb{g} + \frac{\lambda}{m}\,L^2 &= 0
   && \text{putting in $\ddot{\vb{r}} = \vb{g} + (\lambda/m)\vb{r}$ and
      $\vb{r}\cdot\vb{r} = L^2$,}
\end{align*}
so
\begin{equation}
   \lambda = -\frac{m\,\bigl(\lvert\vb{v}\rvert^2 + \vb{r}\cdot\vb{g}
   \bigr)}{L^2} .
   \label{eq:la-multiplier}
\end{equation}
The function \texttt{rod\_multiplier} of \texttt{mdlab} evaluates it. The
rod's force $\lambda\vb{r}$ is $-\lambda L$ times $-\vb{r}/L$, the arrow of
length one from the bob to the pivot, so the rod pulls the bob towards the
pivot with the force $F_{\mathrm{rod}} = -\lambda L$, which is negative
when the rod pushes. With the bob at the angle $\theta$ from the vertical,
$\vb{r} = L(\sin\theta, -\cos\theta)$, so $\vb{r}\cdot\vb{g} =
gL\cos\theta$, and $\lvert\vb{v}\rvert = L\lvert\dot\theta\rvert$
(\cref{sec:la-coordinates}). Then
\begin{equation*}
   F_{\mathrm{rod}} = -\lambda L
   = \frac{m\,\bigl(L^2\dot\theta^2 + gL\cos\theta\bigr)}{L}
   = m\,\bigl(g\cos\theta + L\dot\theta^2\bigr) :
\end{equation*}
the part of the weight along the rod plus the inward pull $mL\dot\theta^2$
that going round a circle needs (\cref{sec:mo-circle}).

Released from rest at the angle $\theta_0$, the bob keeps its energy $K +
U$ of \cref{sec:la-coordinates} at its starting value $mgL(1 -
\cos\theta_0)$, so $\frac12 mL^2\dot\theta^2 = mgL(1 - \cos\theta_0) -
mgL(1 - \cos\theta) = mgL(\cos\theta - \cos\theta_0)$; multiplying by
$2/(mL)$, $L\dot\theta^2 = 2g(\cos\theta - \cos\theta_0)$ and
\begin{equation*}
   F_{\mathrm{rod}} = mg\,(3\cos\theta - 2\cos\theta_0) .
\end{equation*}
\Cref{fig:la-tension} follows a bob of \SI{0.5}{\kilogram} on a rod of
\SI{1.2}{\metre} released from \ang{120}, computed in $x$ and $y$ with the
multiplier: the length stays \SI{1.2}{\metre} to within \SI{1e-11}{\metre},
and the force agrees with the formula at every angle: at the bottom it is
$mg\,(3 - 2\cos\ang{120}) = mg\,(3 + 1) = 4mg = 4\times0.5\times9.80665 =
\SI{19.613}{\newton}$. It is negative where $3\cos\theta < 2\cos\theta_0$,
beyond $\cos\theta = \frac23\cos\theta_0 = -\frac13$, here beyond
\ang{109.47}: the rod pushes the bob outwards. A string, which can only
pull, could not follow this swing at all: at the release point the force
needed is $mg\cos\ang{120} = -\frac12 mg$, a push, so a string is slack
from the start. \Cref{ex:la-slack} shows that a string stays taut only for
swings of at most \ang{90}, the limit of \cref{sec:os-anharmonic}. The
angle \ang{109.47} belongs to this swing of \ang{120}: between \ang{90} and
\ang{109.47} the bob, though above the pivot, moves fast enough that the
rod still pulls.

\begin{figure}[htb]
\centering
\includegraphics{ch06_lagrange/tension}
\caption{The force of the rod on a pendulum's bob of \SI{0.5}{\kilogram}
on a rod of \SI{1.2}{\metre}, released from rest at \ang{120}, divided by
the bob's weight, against the angle over one swing: from the Lagrange
multiplier of the constraint (solid) and from energy, $3\cos\theta - 2\cos
\theta_0$ (dashed), which the solid line hides. Below the dotted line the
rod pushes.}
\label{fig:la-tension}
\end{figure}

\subsection*{Bonds of fixed length}

A bond held at the length $d$ between atoms $i$ and $j$ is the constraint
$\chi = \frac12(\lvert\vb{r}_j - \vb{r}_i\rvert^2 - d^2) = 0$. Its gradient
with respect to $\vb{r}_j$ is $\vb{r}_j - \vb{r}_i$, and with respect to
$\vb{r}_i$ it is $-(\vb{r}_j - \vb{r}_i)$ (\cref{ex:la-bond}), so the
multiplier gives the two atoms equal and opposite forces along the bond, as
the third law requires. Freezing the fastest bonds of a molecule this
way, those to hydrogen, whose periods of \SI{9}{\femto\second} for O--H and
\SI{11}{\femto\second} for C--H limit the step of a simulation (\cref{sec:os-molecules}), lets the steps grow;
Chapter~11 finds the multipliers at every step with the method called
SHAKE.

\begin{keyidea}
To find the largest or smallest value of $f$ subject to $\chi = 0$, solve
$\grad f = \lambda\grad\chi$ with $\chi = 0$. In mechanics a constraint $\chi = 0$
that does no work enters as the force $\lambda\grad\chi$, and $\lambda$,
fixed by differentiating the constraint twice, measures the force needed
to keep it: for a pendulum it gives the force of the rod, $-\lambda L$.
\end{keyidea}

\notebook{06}{the pendulum in $x$ and $y$ with its multiplier}

\begin{selfcheck}
A pendulum is released from rest at \ang{90}. What force does the rod
exert at the bottom of the swing, in units of the bob's weight?
\end{selfcheck}
